\documentclass[10pt,a4paper]{amsart}
\usepackage[margin=19mm]{geometry}
\usepackage{amsmath,amssymb,mathtools}
\usepackage[scr=rsfs,cal=euler]{mathalfa}
\usepackage{xfrac}
\usepackage[unicode,hidelinks,bookmarks=false]{hyperref}
\newcommand{\ie}{i.e.\ }
\newcommand{\cf}{cf.\ }
\newcommand{\resp}{respectively\ }
\newcommand{\Subsection}{\subsection}
\providecommand{\nobreakdash}{\nobreak\hbox{-}}
\setlength{\emergencystretch}{2em}
\multlinegap0pt
\usepackage{bm}
\usepackage{bbm}
\usepackage{dsfont}
\usepackage{xspace}
\usepackage{cancel}
\usepackage{mathtools}
\usepackage{amsthm}
\makeatletter
\@ifundefined{lemma}{\newtheorem{lemma}{Lemma}}{}
\makeatother
\newcommand{\Ccal}{\mathcal{C}} % modal frame space cat param
\newcommand{\Dcal}{\mathcal{D}} % rel space cat param
\newcommand{\Fcal}{\mathcal{F}} % point functor
\newcommand{\Ocal}{\mathcal{O}} % open sets
\newcommand{\Pbb}{\mathbb{P}}
\newcommand{\dirim}{^{\rightarrow\!}}
\newcommand{\interiorop}[1]{{#1}^\circ}
\title[Topological Dualities for Modal Algebras]{Topological Dualities for Modal Algebras}
\author{Matthew Collinson}
\date{Source: arXiv:2604.20603v1 (CC BY 4.0)}
\begin{document}
\maketitle

\noindent\emph{Source and licence.} Topological Dualities for Modal Algebras. Version arXiv:2604.20603v1 (CC BY 4.0), distributed under the Creative Commons Attribution 4.0 International licence, as recorded in the arXiv licence metadata for this exact version and shown on the abstract page https://arxiv.org/abs/2604.20603v1. Full rights notice: licenses/arxiv-2604.20603-CC-BY-4.0.txt.\par\smallskip

\noindent\textit{Editorial scope.} Counted unit: the Lemma whose source label is \texttt{result:sober} in the article (source0.tex lines 1223--1226, source label \texttt{result:sober}), delivered together with the article's own complete original proof of it (source0.tex lines 1228--1296, one \texttt{proof} environment of 69 lines). No mathematics is written, repaired or re-derived: statement and proof are the article's own text line for line. Required context retained uncounted: tripmot (lines 813--823); pointFunctorOnMorphs (lines 1127--1132); psiXp (lines 1145--1148); psiXF (lines 1150--1158); psiXa (lines 1150--1158). Every definition, display and label of those lines is retained unchanged. Omitted from this package and not redistributed: 1--812, 824--1126, 1133--1144, 1149--1149, 1159--1222, 1227--1227, 1297--1736. Nothing inside a retained excerpt is omitted or altered: every retained body is delivered exactly as the article's lines are written, so no omission receipt of any kind accompanies this package. Citations: the retained excerpts carry no citation command at all, so no bibliography entry of the article is transcribed and no reference is redistributed. Reference adaptation: none was needed and none was made; every reference inside the delivered unit resolves inside it, so no unresolved reference or citation is produced. Printed slips kept literal: no printed word, symbol, hypothesis or number of the counted statement or of its proof was altered. Layout and class: the article's own class is \texttt{article} with packages including a4wide, graphicx, pict2e, amsmath, amsfonts, amssymb, amsthm, makecell, mathrsfs, bandd2025; the wrapper is the lane's pinned \texttt{amsart} editorial wrapper at 10pt on A4, which restates the theorem-like environments the retained text needs and adds the article's own macro definitions that the retained text needs (\texttt{\textbackslash Ccal}, \texttt{\textbackslash Dcal}, \texttt{\textbackslash Fcal}, \texttt{\textbackslash Ocal}, \texttt{\textbackslash Pbb}, \texttt{\textbackslash dirim}, \texttt{\textbackslash interiorop}), with extra wrapper packages (bm, bbm, dsfont, xspace, cancel, mathtools). The delivered numbering is the unit's own. Assets: no retained span contains a figure environment, a graphics-inclusion command, a PDF-page inclusion, a table or a TikZ picture; no image file is copied into this package, the package declares no asset dependency and no image file is redistributed with it. Licence: version arXiv:2604.20603v1 of this article is distributed under the Creative Commons Attribution 4.0 International licence, as recorded in the arXiv licence metadata for this exact version and shown on the abstract page https://arxiv.org/abs/2604.20603v1. A full rights notice is delivered at licenses/arxiv-2604.20603-CC-BY-4.0.txt. Characters: every retained excerpt is pure ASCII, so the delivered engine, XeLaTeX with its default Latin Modern text font, reports no missing character.
\begin{lemma}
\label{tripmot}
Let $(p,a,[F]) \in \Pbb _A$. Then for all $c \in A$ 
\begin{enumerate}
\item $c \leq a$ iff $R_A \dirim (p,a,[F])\cap\phi _A (c) 
= \emptyset $ 
\item $c \in  F$ iff 
$R_A \dirim (p,a,F)
\subseteq \phi _A (c)$.
\end{enumerate}
\end{lemma}

\begin{equation}
\label{pointFunctorOnMorphs}
 \mathcal{F} (f)(q,b,[G]) = 
( q \circ f ,                           \bigvee \{ c \in A \mid f(c) \leq b \} ,  
[\{ c \in A \mid f(c) \in G \} ]). 
\end{equation}

\begin{equation}
\label{psiXp}
p_x (U) = 1  \mbox{ iff }  x \in U 
\end{equation}

\begin{align} 
\label{psiXa}
a_x  & =  
\interiorop{(
X \setminus R\dirim(x))
}    \\
\label{psiXF}
F_x  & =  \{ U \in \Ocal \mid R\dirim (x) \subseteq U \} . 
\end{align}

\begin{lemma}
\label{result:sober}
If $A$ is a frame in the category $\Ccal$, then $\mathcal{F} (A)$ is $\Dcal$-sober.
\end{lemma}

\begin{proof}
We give the proof in the case where points are triples. For the proof where points are pairs, elide the filter component of the triples.

We have 
$\mathcal{F}(\phi _A) \psi _{ \mathcal{F}(A)} = 1 _{\mathcal{F}(A)}$ 
by the triangle identities; we will show 
$\psi _{ \mathcal{F}(A) } \mathcal{F} (\phi _A) = 1 _{\mathcal{F}\Omega\mathcal{F}(A)}$.
Let $(p, U , F ) \in \mathcal{F} 
\Omega \mathcal{F} (A)$.

Let
\begin{equation*}
x = \mathcal{F}(\phi _A)(p, U, F ) 
 =  (p \circ \phi _A, e , K ) 
\end{equation*}
where 
\begin{align*}
    e &= \bigvee \{ c \in A | \phi _A (c) \subseteq U\} \\
    K &= \{ c \in A | \phi _A (c) \in F \} 
\end{align*}
as in (\ref{pointFunctorOnMorphs}).
Let $\psi _{\mathcal{F} A} (x) = (p_x , a_x , F_x)$ as in (\ref{psiXp}), (\ref{psiXa}), (\ref{psiXF}), where now $X = \Fcal(A)$.
%\ref{definition:fhashTriple}, 
%i.e.
%\begin{equation*}
%\begin{array}{c}
%p_x (U) = 1 ~ \Leftrightarrow ~ x \in U ~                          \\ 
%a_x ^i = \interiorop{(\{ y \in \mathcal{F} A\mid x\not\!\! R_A ^i y \})} \\
%F_x ^i = \{ U\in\Omega\mathcal{F} A\mid R_A ^i (x) \subseteq U \}  . 
%\end{array}
%\end{equation*}
We show $(p,U,F) = (p_x , a_x , F_x)$. %The proof for the first two components is as in HILKEN. The proof for the third component is as in THESIS. 

For any $c \in A$, we have $p_x ( \phi _A (c))=1$ iff $x \in \phi _A(c)$ iff 
$p ( \phi _A (c))=1$, so $p_x = p$ as $\phi_A$ is surjective.

For any $c \in A$, it is the case that $\phi _A (c) \subseteq U$ iff 
$c \leq e$, by definition of $e$ and because $\phi_A$ is a frame morphism. 
However, 
%\begin{equation*}
%\begin{array}{rl}
$\phi _A(c) \subseteq a_x$ 
iff (by definition of $a_x$)
$\phi _A (c) \subseteq  X \setminus R_A\dirim (x)$ 
%\\
%\Leftrightarrow & \forall (q,b,G) \in \mathbb{P} _A . x R_A (q,b,G)  
%                  \Rightarrow q(a)=0        %                         \\ 
%\Leftrightarrow & a \leq e                  %                      .
%\end{array}
%\end{equation*}
iff $c \leq e$, by definition of $x$ and Lemma~\ref{tripmot}.
Therefore $U = a_x$. 

For all $c \in A$ we have 
%\begin{align*}
$\phi _A (c) \in F_x$ iff   
%& \Longleftrightarrow 
$R_A \dirim (x)  \subseteq \phi_A (c)$
%\\ 
iff 
%& \Longleftrightarrow  
$c \in K$ 
%\\
iff 
%&\Longleftrightarrow  
 $\phi _A (c) \in F$,                             
%\end{align*}
using Lemma~\ref{tripmot}. Therefore $F_x  = F$.
\end{proof}

\end{document}
